% ============================================================================
\section{Prueba de usabilidad}\label{sec:a5-usabilidad}

La prueba se aplicó el jueves 20 de agosto de 2026 sobre el producto desplegado, con cinco
participantes reales y dos bloques complementarios. Este apartado documenta el estudio completo:
objetivos, participantes, escenarios, métricas, materiales, procedimiento, ejecución, análisis,
hallazgos y recomendaciones. La Tabla \ref{tab:usab-diez-pasos} señala dónde queda cubierto cada
uno de los diez pasos que la instrucción de la actividad propone para desarrollar una prueba de
usabilidad, de manera que cada paso pueda verificarse por separado.

\begin{uxtablalarga}{C{0.7cm}L{3.6cm}L{3.2cm}Y}{Los diez pasos de la instrucción y el apartado que cubre cada uno\label{tab:usab-diez-pasos}}{\thd{Paso} & \thd{Qué pide la instrucción} & \thd{Dónde queda cubierto} & \thd{Evidencia principal}}
1 & Definir los objetivos de la evaluación & Apartado \ref{subsec:usab-objetivos}, Objetivos & Cinco objetivos declarados: efectividad, eficiencia, comprensión, facilidad y satisfacción \\
2 & Identificar a los usuarios objetivo, con su experiencia y las características relevantes & Apartado \ref{subsec:usab-participantes}, Participantes y reclutamiento & Tabla de los cinco participantes con perfil, experiencia y tratamiento en el reporte \\
3 & Desarrollar escenarios de prueba realistas, con tareas claras & Apartado \ref{subsec:usab-diseno}, Diseño de la prueba y escenarios & Seis tareas estandarizadas y tres recorridos de trabajo, cada uno con su criterio de éxito \\
4 & Elegir las métricas de usabilidad en función de los objetivos & Apartado \ref{subsec:usab-materiales} y el apartado ``Descripción de métricas UX'' & Tiempo por tarea, tasa de éxito, errores, ayudas, SEQ y SUS, cada métrica con su meta \\
5 & Crear los materiales de prueba: prototipo, wireframe, producto o sistema real & Apartado \ref{subsec:usab-materiales}, Materiales & El producto real desplegado, con la versión fijada y datos sintéticos de semilla conocida \\
6 & Reclutar participantes que coincidan con el grupo de usuarios objetivo & Apartado \ref{subsec:usab-participantes}, Participantes y reclutamiento & Muestreo intencional de cinco profesionales que usan datos en su trabajo, con consentimiento \\
7 & Realizar la prueba con instrucciones y escenarios, recogiendo lo que la persona expresa & Apartado \ref{subsec:usab-materiales}, Procedimiento & Protocolo de siete pasos aplicado el 20 de agosto de 2026, una sesión presencial y dos remotas \\
8 & Recopilar datos cuantitativos y cualitativos & Apartados \ref{subsec:usab-resultados-a}, \ref{subsec:usab-resultados-b} y \ref{subsec:usab-sus} & Tiempo, éxito, errores y SEQ por tarea; finalización, duración y facilidad por recorrido; SUS de los cinco \\
9 & Analizar e interpretar los resultados frente a las métricas definidas & Apartados \ref{subsec:usab-sus} y \ref{subsec:usab-hallazgos} & Síntesis contra metas, media SUS de 89.0 con su intervalo de confianza y hallazgos por frecuencia y tipo \\
10 & Informar hallazgos y recomendaciones con ejemplos concretos y ayudas visuales & Apartados \ref{subsec:usab-hallazgos}, \ref{subsec:usab-recomendaciones} y \ref{subsec:usab-limitaciones} & H1 a H8 con su evidencia, R1 a R7 priorizadas y las limitaciones declaradas del estudio \\
\end{uxtablalarga}

\subsection{Objetivos}\label{subsec:usab-objetivos}

La prueba evaluó si profesionales habituados a consultar información podían completar recorridos
principales, comprender fuentes y estados, y reconocer oportunidades de mejora sin conocer
previamente la estructura del portal. Los objetivos fueron:

\begin{uxlist}
  \lead{Efectividad}{comprobar que las tareas o actividades asignadas terminaran con el resultado correcto.}
  \lead{Eficiencia}{medir tiempo, errores y solicitudes de ayuda cuando el protocolo lo permitiera.}
  \lead{Comprensión}{verificar que fuente, estado, proyección, permisos y opciones de extracción se interpretaran correctamente.}
  \lead{Facilidad}{registrar facilidad percibida después de cada tarea o al terminar el recorrido.}
  \lead{Satisfacción}{medir la percepción global mediante la escala SUS.}
\end{uxlist}

\subsection{Participantes y reclutamiento}\label{subsec:usab-participantes}

Se utilizó un muestreo intencional de cinco profesionales con experiencia real en consulta,
validación o análisis de información. Los criterios de inclusión fueron utilizar datos en el
trabajo, realizar búsquedas, cruces o validaciones de fuentes y poder comparar el prototipo con un
proceso habitual. Todos otorgaron su consentimiento. Dos autorizaron publicar su nombre y tres
solicitaron anonimato.

\begin{uxtablalarga}{C{1.0cm}L{3.8cm}Y L{2.4cm}}{Participantes de la prueba de usabilidad}{\thd{Código} & \thd{Participante o área} & \thd{Experiencia relevante} & \thd{Tratamiento en el reporte}}
C & Carlos René Flores M. & Científico de datos con experiencia en BBVA AutoMarket; utiliza modelos, datos y herramientas de extracción & Nombre autorizado \\
V & Verónica Itzel Flores González & Analista en la Dirección de Información del Sistema Financiero; valida cifras y prepara información & Nombre autorizado \\
A1 & Oficina de Información de Crédito al Consumo y Comisiones & Analista de datos, aproximadamente 12 años de experiencia & Anónimo \\
A2 & Subgerencia de Análisis Preventivo & Consulta analítica y directiva, aproximadamente 15 años de experiencia & Anónimo \\
A3 & Subgerencia de Transparencia de Servicios Financieros & Consulta de información operativa, aproximadamente 6 años de experiencia & Anónimo \\
\end{uxtablalarga}

Los participantes C y V realizaron un protocolo estandarizado de seis tareas. A1, A2 y A3
realizaron recorridos vinculados con su trabajo habitual; estas sesiones ocurrieron el jueves 20
de agosto de 2026, en computadora y con la misma versión del prototipo. A1 participó de forma
presencial y A2 y A3 de forma remota.

\textbf{Tratamiento de los datos personales.} Solo C y V autorizaron de forma expresa que su
nombre aparezca en el reporte. A1, A2 y A3 se identifican únicamente por código: de ellos no se
publica nombre, correo, teléfono ni ningún otro identificador, y su columna conserva solo el dato
de contexto que da sentido a su recorrido. Los consentimientos y las notas de sesión se conservan
fuera de este documento.

\subsection{Diseño de la prueba y escenarios}\label{subsec:usab-diseno}

La evaluación integró dos bloques complementarios. El primero produjo métricas comparables por
tarea sobre las seis interfaces. El segundo permitió comprobar relevancia y descubribilidad en
actividades cercanas al contexto laboral de tres perfiles adicionales. Sus resultados no se
fusionan cuando la unidad de medición es distinta.

\textbf{Por qué dos bloques.} La tarea guiada mide y el escenario abierto revela para qué se
usaría el producto: una sola de las dos formas deja fuera la mitad de la evidencia. La evaluación
de sistemas conversacionales con participantes reales en escenarios controlados sostiene esa
combinación y da la forma de la tarea del asistente, T5, donde lo que se juzga no es solo si la
respuesta llega, sino si la persona reconoce qué consulta y qué fuente la respaldan
(Zhu et al., 2026).

\textbf{Una prueba por interfaz, ligada a su funcionalidad principal.} La instrucción de la
actividad pide, en consecuencia, al menos una prueba de usabilidad relacionada con la
funcionalidad principal de cada interfaz diseñada. La Tabla \ref{tab:usab-interfaz-tarea} responde
a esa exigencia punto por punto: recorre las siete interfaces publicadas en la Actividad 4, nombra
la función que cada una debe resolver, la prueba que la ejercitó y la métrica con la que se juzgó.
Cinco interfaces se ejercitaron con una tarea cronometrada propia. Inicio y Gobierno del dato no
tuvieron tarea propia y se observaron dentro de los recorridos del bloque B, así que su renglón
reporta lo que se observó y deja dicho que no hubo medición cronometrada, en lugar de presentar
una cifra que no se tomó.

\begin{uxtablalarga}{L{2.7cm}L{3.3cm}L{2.5cm}Y}{Correspondencia entre interfaz diseñada, funcionalidad principal, prueba y métrica\label{tab:usab-interfaz-tarea}}{\thd{Interfaz de la Actividad 4} & \thd{Funcionalidad principal} & \thd{Prueba que la ejercita} & \thd{Métrica y resultado observado}}
Acceso por perfil & Autenticar y llevar a cada rol al espacio que le corresponde & T1, bloque A & Éxito y tiempo: 2 de 2 con éxito, en 24 y 29 s, bajo la meta de 45 s \\
Inicio & Ofrecer el punto de partida de la búsqueda y los accesos propios del rol & Arranque de los tres recorridos del bloque B & Finalización sin ayuda: 3 de 3; H2 registra la duda sobre dónde empezar. Sin tarea cronometrada propia \\
Exploración y extracción, con su rama de tableros e indicadores & Localizar un campo por concepto y leer la serie que lo acompaña & T2 y T4, bloque A & Tiempo y metadatos correctos: T2 en 51 y 62 s con un error por persona; T4 dio un éxito y un éxito parcial, con SEQ 5 y 4 \\
Gobierno del dato & Publicar fuente, responsable y certificación de cada campo & Lectura de la ficha dentro de T2 y de los recorridos del bloque B & Reconocimiento de fuente y certificación sin ayuda: 3 de 3 en el bloque B (H6). Sin tarea cronometrada propia \\
Asistente conversacional & Responder en lenguaje natural mostrando la consulta y la fuente que respaldan la cifra & T5, bloque A & Fuente identificada y ayudas: 2 de 2 sin ayuda, en 49 y 57 s, con SEQ 6 y 6 \\
Administración & Revisar una cuenta y su acceso sin modificarla & T6, bloque A & Éxito y tiempo: 2 de 2 con éxito, en 58 y 71 s, con un error registrado en V \\
Exportación & Solicitar un archivo reutilizable y seguir su estado y su vigencia & T3, bloque A & Tiempo y comprensión del estado: 2 de 2 con éxito, en 64 y 79 s, bajo la meta de 120 s \\
\end{uxtablalarga}

\subsubsection{Bloque A. Seis tareas estandarizadas}

\begin{uxtablalarga}{C{0.8cm}L{2.5cm}Y L{4.0cm}}{Tareas estandarizadas para C y V}{\thd{T} & \thd{Interfaz} & \thd{Escenario de prueba} & \thd{Indicadores y meta}}
1 & Acceso por perfil & Entrar como analista y confirmar que se alcanzó el espacio correspondiente & Éxito 100\%; tiempo $\leq45$ s; cero selecciones equivocadas \\
2 & Catálogo & Encontrar el campo de cobertura de liquidez e identificar fuente, responsable y certificación & Éxito 100\%; tiempo $\leq90$ s; máximo un error; tres metadatos correctos \\
3 & Exportaciones & Solicitar la cartera de crédito en XLSX y localizar el estado y la vigencia del enlace & Éxito 100\%; tiempo $\leq120$ s; máximo un error; estado y caducidad comprendidos \\
4 & Tablero & Explicar la tendencia de cobertura de liquidez y distinguir dato observado de proyección & Interpretación correcta; tiempo $\leq120$ s; SEQ $\geq5$ \\
5 & Asistente & Preguntar por liquidez e identificar qué consulta y fuente respaldan la respuesta & Éxito 100\%; tiempo $\leq90$ s; fuente identificada; cero ayudas \\
6 & Administración & Cambiar al perfil administrador, localizar una cuenta y explicar cómo cambiar o retirar su acceso sin modificarla & Éxito 100\%; tiempo $\leq90$ s; cero acciones accidentales \\
\end{uxtablalarga}

\subsubsection{Bloque B. Recorridos vinculados con el trabajo}

\begin{uxtablalarga}{C{1.0cm}L{3.0cm}Y L{3.5cm}}{Actividades observadas para A1, A2 y A3}{\thd{Código} & \thd{Perfil utilizado} & \thd{Escenario de evaluación} & \thd{Criterio de éxito}}
A1 & Analista & Localizar fuentes y campos de interés, revisar cómo se exportan los datos y consultar indicadores & Completar el recorrido y reconocer la relación entre campos y fuentes sin ayuda \\
A2 & Analista y directivo & Encontrar nuevas fuentes, revisar alertas, indicadores y series, y continuar hacia la extracción & Completar las actividades de ambos perfiles y localizar la información requerida sin ayuda \\
A3 & Consulta operativa & Encontrar fuentes útiles para información operativa y revisar la información disponible sobre ellas & Localizar la fuente y comprender la información presentada sin ayuda \\
\end{uxtablalarga}

\subsubsection{Cobertura realmente observada}

La matriz distingue disponibilidad del prototipo de uso directo. Una explicación posterior no se
contabiliza como tarea observada.

\begin{uxtabla}{L{3.0cm}L{4.0cm}Y}{Cobertura de interfaces por bloque}
  \uxheadrow \thd{Interfaz} & \thd{Bloque A, estandarizado} & \thd{Bloque B, por contexto} \\
  Acceso por perfil & C y V & No evaluado directamente \\
  Catálogo y fuentes & C y V & A1, A2 y A3 \\
  Exportaciones & C y V & A1 y A2 \\
  Tableros e indicadores & C y V & A1 y A2; A3 recibió explicación posterior \\
  Asistente & C y V & No evaluado directamente \\
  Administración & C y V & No evaluado directamente \\
\end{uxtabla}

\subsection{Materiales, métricas y procedimiento}\label{subsec:usab-materiales}

\textbf{El material de prueba fue el producto real desplegado.} La instrucción admite prototipos,
wireframes o el producto o sistema real, y esta prueba usó lo último: las cinco personas
trabajaron sobre la versión del portal publicada en la nube, la misma que reúne las siete
pantallas de la Actividad 4 con el tema institucional, sin maquetas intermedias ni pantallas
dibujadas para la ocasión. La versión quedó fijada antes de la primera sesión y no se modificó
entre una sesión y otra, de modo que los cinco resultados hablan del mismo producto. El acceso
ocurrió por la dirección pública del despliegue, sin instalación local ni descarga de archivos:
\urlDemoWeb. Se trata de un prototipo de alta fidelidad con datos sintéticos, no conectado a
sistemas reales de ninguna institución.

\textbf{Entorno de ejecución.} La interfaz web corre en un servicio administrado de Google Cloud
Run, la API en un servicio privado del mismo entorno y la persistencia en una base de datos
PostgreSQL administrada. El contenido que los participantes consultaron proviene de los silos
sintéticos que genera la semilla fija 20260720, así que cualquier persona que repita el
procedimiento ve las mismas cifras. Cada participante entró con una de las cuentas de
demostración por perfil, en computadora y por navegador. La verificación del despliegue del 22 de
agosto de 2026 midió 2.4 s en el arranque en frío de la primera visita y 0.17 s en las visitas
siguientes, y confirmó que las rutas protegidas redirigen al acceso cuando no hay sesión: la
espera que encuentra quien abre la dirección por primera vez está medida, no supuesta.

\textbf{Métricas elegidas y métricas descartadas.} Las métricas se eligieron por objetivo,
siguiendo la clasificación de Albert y Tullis (2013). Del desempeño se tomaron tiempo por tarea,
tasa de éxito, errores y solicitudes de ayuda (cap. 4); de las auto-reportadas, el SEQ después de
cada tarea y el SUS al cierre (cap. 6); y la lectura conjunta de ambas familias es la que permite
afirmar que una tarea se completó pero costó, como ocurre en el tablero (cap. 8). Las métricas
conductuales y fisiológicas del capítulo 7 ---seguimiento ocular, expresión facial, respuesta
galvánica--- se descartaron con causa: exigen instrumentación de laboratorio que dos sesiones
remotas no admiten y ninguno de los cinco objetivos las necesita. La retención y la tasa de
conversión, que la instrucción menciona entre las métricas comunes, quedaron fuera porque
requieren un producto en operación con uso repetido y no una sesión única. La accesibilidad se
midió fuera de esta prueba, en la verificación de contraste de la Actividad 4, porque su
instrumento es la medición del tema y no la observación de una tarea.

\textbf{Del evaluador prototipo a la persona real.} Esta prueba es el último eslabón de una cadena
que empezó sintética. En la Actividad 3, ocho evaluadores prototipo condicionados por las ocho
personas de la Actividad 1 resolvieron el card sorting y la prueba de árbol bajo el protocolo de
PerceptUI (Bougie et al., 2026); en la Actividad 4, esos mismos ocho evaluadores recorrieron las
siete capturas del prototipo antes de que ninguna persona lo viera. La literatura de 2026 encuadra
esa secuencia: la evaluación automatizada de UX con interacción simulada aplica SUS, SEQ y
pensamiento en voz alta sintéticos como complemento previo, nunca como sustituto de la evidencia
humana (Tan et al., 2026). El proyecto ejecutó exactamente esa secuencia y la
terminó en personas, con una separación deliberada: las escalas SUS y SEQ se aplicaron únicamente
a los cinco participantes reales y ninguna cifra de origen sintético se promedia con las suyas.
La ronda barata sirvió para llegar a la sesión con personas sin gastarla en defectos que se ven a
simple vista.

\textbf{Procedimiento.} Cada sesión siguió estos siete pasos:

\begin{uxpreg}
  \item Presentar propósito, consentimiento y aviso de datos sintéticos.
  \item Registrar experiencia y asignar un código cuando se solicitó anonimato.
  \item Abrir el prototipo en computadora sin anticipar la ubicación de funciones.
  \item En el bloque A, leer una tarea, cronometrar, registrar errores y aplicar SEQ de 1 a 7.
  \item En el bloque B, observar recorridos laborales, duración total, dudas y solicitudes de ayuda.
  \item Aplicar una valoración general de facilidad de 1 a 7 al bloque B.
  \item Aplicar SUS al terminar y separar lo usado de lo conocido únicamente por explicación.
\end{uxpreg}

El bloque A registró tiempo, éxito, errores, ayudas y SEQ por tarea. El bloque B registró
finalización, solicitudes de ayuda, duración total y facilidad general; no se reconstruyeron
tiempos ni errores por tarea que no quedaron anotados. Los cinco resultados SUS sí comparten la
misma escala y se presentan juntos. Una de las tres respuestas anónimas se recuperó posteriormente
y se considera retrospectiva.

Junto a esas medidas se conservó el registro cualitativo de la observación, del que salen los
hallazgos H1 a H4 y H7: ninguna cifra de tiempo o de éxito los habría revelado por sí sola.

\subsection{Resultados del bloque estandarizado}\label{subsec:usab-resultados-a}

Los tiempos están en segundos; E significa éxito y P, éxito parcial.

\begin{uxtablalarga}{C{0.7cm}L{2.2cm}C{1.2cm}C{1.0cm}C{1.0cm}C{1.2cm}C{1.0cm}C{1.0cm}Y}{Resultados de C y V por tarea}{\thd{T} & \thd{Interfaz} & \thd{C tiempo} & \thd{C éxito} & \thd{C err.} & \thd{V tiempo} & \thd{V éxito} & \thd{V err.} & \thd{SEQ C/V}}
1 & Acceso & 24 & E & 0 & 29 & E & 0 & 7 / 7 \\
2 & Catálogo & 51 & E & 1 & 62 & E & 1 & 6 / 6 \\
3 & Exportación & 64 & E & 1 & 79 & E & 1 & 6 / 5 \\
4 & Tablero & 76 & E & 0 & 94 & P & 2 & 5 / 4 \\
5 & Asistente & 49 & E & 0 & 57 & E & 0 & 6 / 6 \\
6 & Administración & 58 & E & 0 & 71 & E & 1 & 6 / 5 \\
\end{uxtablalarga}

\begin{uxtabla}{L{3.3cm}C{2.6cm}C{2.4cm}Y}{Síntesis del bloque estandarizado}
  \uxheadrow \thd{Indicador} & \thd{Resultado} & \thd{Meta} & \thd{Lectura} \\
  Tasa de éxito & 11 de 12, 91.7\% & $\geq90\%$ & Cumple; el tablero concentra el éxito parcial \\
  Tiempo medio por tarea & 59.5 s & Informativa & Ninguna media por interfaz supera su umbral \\
  Errores observados & 7 & Máximo 1 por tarea y persona & V supera el umbral únicamente en tablero \\
  Ayudas de moderación & 0 & Máximo 1 por persona & Cumple \\
  SEQ promedio & 5.8 de 7 & $\geq5$ & Cumple; tablero obtiene 4.5 \\
\end{uxtabla}

\subsection{Resultados de los recorridos por contexto}\label{subsec:usab-resultados-b}

\begin{uxtabla}{C{1.0cm}C{2.2cm}C{2.1cm}C{2.0cm}C{1.8cm}Y}{Resultados de A1, A2 y A3}
  \uxheadrow \thd{Código} & \thd{Finalización} & \thd{Ayuda} & \thd{Duración} & \thd{Facilidad} & \thd{Observación principal} \\
  A1 & Completa & Ninguna & 30--45 min & 5 de 7 & Dudó qué parte utilizar primero; los campos ayudaron a reconocer fuentes \\
  A2 & Completa & Ninguna & 45 min & 6 de 7 & Encontró fuentes e indicadores; solicitó agregaciones o transformaciones sencillas \\
  A3 & Completa & Ninguna & 40 min & 7 de 7 & El recorrido de fuentes fue claro; expresó interés posterior en alertas y tableros \\
\end{uxtabla}

Las tres personas completaron sus actividades sin ayuda. La facilidad media fue 6.0 de 7. Las
sesiones duraron entre 30 y 45 minutos; el intervalo de A1 impide calcular una media temporal
exacta. A3 no utilizó directamente alertas o tableros y su interés no se contabilizó como éxito.

\subsection{Resultados SUS de los cinco participantes}\label{subsec:usab-sus}

\begin{uxtabla}{C{1.1cm}Y C{2.0cm}}{Puntuaciones SUS consolidadas}
  \uxheadrow \thd{Código} & \thd{Registro disponible} & \thd{SUS} \\
  C & Puntuación final del protocolo estandarizado & 87.5 \\
  V & Puntuación final del protocolo estandarizado & 82.5 \\
  A1 & 4, 2, 4, 3, 4, 1, 5, 2, 5, 1 & 82.5 \\
  A2 & 4, 1, 4, 2, 5, 1, 5, 1, 5, 1 & 92.5 \\
  A3 & 5, 1, 5, 1, 5, 1, 5, 1, 5, 1 & 100.0 \\
\end{uxtabla}

El promedio combinado fue \textbf{89.0 de 100}, superior a la meta de 75. El registro consolidado
de C y V conserva la puntuación final; las notas actuales conservan además las respuestas por
reactivo de A1, A2 y A3. La puntuación muestra una percepción favorable en esta muestra, pero no
elimina las fricciones observadas ni se generaliza estadísticamente.

\clearpage
\subsection{Hallazgos}\label{subsec:usab-hallazgos}

\begin{uxtablalarga}{C{1.0cm}L{3.0cm}Y C{1.4cm}L{1.8cm}}{Hallazgos consolidados}{\thd{ID} & \thd{Hallazgo} & \thd{Evidencia e impacto} & \thd{Frec.} & \thd{Tipo}}
H1 & Proyección y dato observado compiten visualmente & V interpretó primero la proyección como último dato; T4 fue parcial y obtuvo SEQ 4 & 1 de 5 & Alta \\
H2 & El inicio de la búsqueda no siempre es evidente & C y V probaron primero el buscador de cabecera; A1 dudó qué parte utilizar primero & 3 de 5 & Media \\
H3 & El selector de exportación comunica tarde sus límites & C y V revisaron las restricciones después de elegir XLSX, aunque completaron la tarea & 2 de 5 & Media \\
H4 & El cambio de perfil requiere confirmar el nuevo contexto & V buscó Administración antes de cambiar de perfil y revisó nuevamente la cabecera & 1 de 5 & Media \\
H5 & La trazabilidad del asistente se reconoce con facilidad & C y V identificaron consulta y fuente sin ayuda y calificaron T5 con SEQ 6 & 2 de 2 expuestos & Fortaleza \\
H6 & Campos y metadatos apoyan el reconocimiento de fuentes & A1 consideró útiles los campos; A2 y A3 localizaron la información requerida con facilidad & 3 de 3 expuestos & Fortaleza \\
H7 & El flujo analítico termina antes de una transformación ligera & A2 solicitó agregaciones o transformaciones sencillas antes de extraer & 1 de 5 & Oportunidad \\
H8 & Alertas e indicadores tienen valor para seguimiento & A2 los utilizó desde el perfil directivo; A3 expresó interés después de conocer su propósito & 1 uso + 1 interés & Señal cualitativa \\
\end{uxtablalarga}

\subsection{Recomendaciones}\label{subsec:usab-recomendaciones}

\begin{uxlist}
  \lead{R1. Separar observado y proyectado}{usar rótulos persistentes, fecha de corte y lenguaje inequívoco para la proyección. Atiende H1.}
  \lead{R2. Unificar la intención de búsqueda}{diferenciar búsqueda global y contextual, conservar el término y describir qué resuelve cada control. Atiende H2.}
  \lead{R3. Llevar límites al selector}{mostrar tamaño, compatibilidad y uso recomendado junto a cada formato. Atiende H3.}
  \lead{R4. Confirmar el contexto al cambiar perfil}{redirigir al inicio del perfil o mostrar un aviso con acceso directo a su espacio. Atiende H4.}
  \lead{R5. Conservar metadatos y trazabilidad visibles}{mantener fuente, responsable, consulta y estado porque sostienen H5 y H6.}
  \lead{R6. Evaluar transformaciones ligeras}{probar con más analistas qué agregaciones aportan valor antes de ampliar el alcance. Responde a H7.}
  \lead{R7. Retestar después de los cambios}{repetir las tareas T2 y T4, cronometrar cada actividad del bloque por contexto, ampliar la muestra hacia perfiles principiantes y comparar contra esta línea base.}
\end{uxlist}

La prioridad es R1 por el impacto de interpretar una cifra directiva de forma incorrecta. R2
aparece en tres participantes y es la fricción más frecuente. R3 y R4 mejoran eficiencia; R5
protege fortalezas. R6 requiere validación antes de convertirse en desarrollo.

\subsection{Limitaciones}\label{subsec:usab-limitaciones}

Los cinco participantes tenían experiencia con información, por lo que la muestra no representa a
personas principiantes. Los dos bloques utilizaron unidades distintas: tiempo y SEQ por tarea en el
estandarizado, y duración y facilidad global en los recorridos laborales. Solo SUS se consolidó.
Las respuestas por reactivo de C y V no están transcritas en las notas actuales, aunque sus
puntuaciones finales sí. Una respuesta anónima se obtuvo retrospectivamente. Finalmente, A3 no
utilizó directamente alertas y tableros. Estas restricciones acotan las afirmaciones, pero dejan
una línea base auditable de cinco participantes reales.

\textbf{Sobre el tamaño de la muestra.} Cinco participantes son pocos y conviene decir con qué
propósito se eligieron. El estudio es de descubrimiento de problemas, no de comparación entre
diseños ni de estimación de un valor poblacional: en esa clase de estudio la muestra se dimensiona
por su capacidad de sacar a la luz los defectos frecuentes, no por la precisión de un promedio
(Albert y Tullis, 2013, cap. 3). El caso guiado que la actividad ofrece como referencia sugiere
entre ocho y diez participantes para obtener significancia estadística (Ramírez Mejía, 2023), y
esta prueba se quedó por debajo de esa cifra. La distancia se responde con tres elementos, en
lugar de omitirla. Primero, el intervalo de confianza del 95 \% para la media SUS es
[79.8, 98.2] ---su cálculo se detalla en el apartado ``Descripción de métricas UX''--- y su
límite inferior queda por encima de la meta de 75 y del promedio de referencia de 68 puntos que la
literatura de SUS utiliza como línea media (Sauro y Lewis, 2016), de modo que la lectura ``la
percepción es favorable'' sobrevive incluso al escenario pesimista de esta muestra. Segundo,
ningún hallazgo se expresa como porcentaje de población: H1 a H8 llevan su frecuencia en absolutos
sobre las personas expuestas a cada interfaz, que es la forma honesta de reportar con esta n.
Tercero, la recomendación R7 fija el retest posterior a los cambios con una muestra ampliada hacia
los perfiles principiantes que este panel no cubre; el ciclo de medir, corregir y volver a medir
es el que sostiene la mejora documentada en herramientas de analítica de autoservicio comparables
(Joarder et al., 2026).
